\documentclass[12pt,a4paper]{article}
\usepackage[czech]{babel}
\usepackage[T1]{fontenc}
\usepackage[utf8]{inputenc}
\usepackage{amsmath,amssymb}
\usepackage{geometry}
\geometry{margin=2.5cm}

\begin{document}

\begin{center}
\Large
\textbf{NMAG~111 -- Lineární algebra I}\\
\vspace{2mm}
\textbf{Vzorový midterm (kapitoly 2--5)}
\normalsize

\vspace{2mm}
Doba řešení: 90 minut \\
Celkový počet bodů: 50
\end{center}

\vspace{4mm}

\textbf{Instrukce.}
\begin{itemize}
  \item Odpovědi pečlivě zdůvodňujte, není-li řečeno jinak.
  \item Nejsou povoleny žádné elektronické pomůcky ani kalkulačka.
  \item Pište čitelně a strukturovaně.
\end{itemize}

\vspace{4mm}

\begin{center}
\begin{tabular}{|c|c|}
\hline
Příklad & Body \\
\hline
1 & 6 \\
2 & 6 \\
3 & 6 \\
4 & 6 \\
5 & 8 \\
6 & 8 \\
7 & 10 \\
\hline
Celkem & 50 \\
\hline
\end{tabular}
\end{center}

\newpage

\textbf{(1) [6 bodů]} Zakroužkujte správnou odpověď, nezdůvodňujte.
Pro získání bodů je nutné odpovědět správně \emph{všechny tři} podotázky.

\medskip
\textbf{ANO / NE}

\begin{enumerate}
\item Nechť $V$ je vektorový prostor nad tělesem $\mathbb{R}$ a $v_1,v_2\in V$.
\begin{enumerate}
  \item Je-li množina $\{v_1,v_2\}$ lineárně závislá, pak aspoň jeden z vektorů $v_1,v_2$ je nulový.
  \item Je-li $v_1\neq 0$ a $v_2 = 3v_1$, pak $\{v_1,v_2\}$ je báze lineárního obalu $\{v_1,v_2\}$.
  \item Každé dva různé vektory ve $V$ tvoří lineárně nezávislou množinu.
\end{enumerate}

\item Nechť $A$ je matice typu $5\times 3$ nad $\mathbb{R}$ hodnosti $3$.
\begin{enumerate}
  \item Jádro lineárního zobrazení určeného maticí $A$ je triviální.
  \item Sloupce matice $A$ tvoří bázi $\mathbb{R}^5$.
  \item Existuje matice $B$ typu $3\times 5$ taková, že $BA=I_3$.
\end{enumerate}

\item Nechť $f\colon V\to W$ je lineární zobrazení.
\begin{enumerate}
  \item Platí $f(0_V)=0_W$.
  \item Je-li $f$ prosté, pak $\ker f=\{0\}$.
  \item Je-li $f$ surjektivní, pak je prosté.
\end{enumerate}
\end{enumerate}

\newpage

\textbf{(2) [6 bodů]} Uveďte definice následujících pojmů.
Pište celými větami, nikoliv pouze schematicky.

\begin{enumerate}
  \item Vektorový prostor nad tělesem.
  \item Lineární kombinace vektorů.
  \item Jádro lineárního zobrazení.
\end{enumerate}

\newpage

\textbf{(3) [6 bodů]} V této úloze není třeba řešení zdůvodňovat, stačí správný výsledek.

\begin{enumerate}
\item Najděte dimenzi a jednu bázi lineárního obalu množiny
\[
\{(1,2,1),(2,4,2),(1,0,1)\}\subset \mathbb{R}^3.
\]

\item Najděte všechna řešení homogenní soustavy lineárních rovnic
\[
\begin{aligned}
x_1 + x_2 - x_3 &= 0,\\
2x_1 + 2x_2 - 2x_3 &= 0.
\end{aligned}
\]

\item Určete hodnost matice
\[
A=
\begin{pmatrix}
1 & 2 & 3\\
2 & 4 & 6\\
1 & 1 & 1
\end{pmatrix}.
\]
\end{enumerate}

\newpage

\textbf{(4) [6 bodů]} Proveďte podrobný výpočet se stručným komentářem.

Najděte bázi jádra a bázi obrazu lineárního zobrazení
$f\colon \mathbb{R}^4\to\mathbb{R}^3$
určeného maticí
\[
A=
\begin{pmatrix}
1 & 0 & 1 & 2\\
0 & 1 & -1 & 1\\
1 & 1 & 0 & 3
\end{pmatrix}.
\]

\newpage

\textbf{(5) [8 bodů]} Následující tvrzení pouze přesně formulujte, \emph{nedokazujte}.
Uveďte všechny předpoklady a vysvětlete použitá označení.

\begin{enumerate}
  \item Tvrzení o vztahu dimenze jádra, dimenze obrazu a dimenze definičního oboru lineárního zobrazení.
  \item Tvrzení o tom, kdy jsou sloupce matice lineárně nezávislé.
\end{enumerate}

\newpage

\textbf{(6) [8 bodů]} Dokažte následující tvrzení. V důkazu se opírejte pouze
o definice a dříve dokázaná základní tvrzení.

\begin{enumerate}
\item Nechť $V$ je vektorový prostor a $v_1,\dots,v_n\in V$.
Dokažte, že pokud je tato množina lineárně nezávislá, pak každý vektor
$v\in \mathrm{span}(v_1,\dots,v_n)$ má právě jednu reprezentaci jako jejich
lineární kombinace.

\item Nechť $f\colon V\to W$ je lineární zobrazení mezi konečně rozměrnými
vektorovými prostory stejné dimenze. Dokažte, že je-li $f$ prosté, pak je
surjektivní.
\end{enumerate}

\newpage

\textbf{(7) [10 bodů]} Odpovědi pečlivě zdůvodněte.

\begin{enumerate}
\item Existuje lineární zobrazení $f\colon \mathbb{R}^4\to\mathbb{R}^3$, které je
\begin{enumerate}
  \item surjektivní?
  \item prosté?
\end{enumerate}
V obou případech svou odpověď zdůvodněte.

\item Najděte všechny matice $A$ typu $2\times 2$ nad $\mathbb{R}$, které splňují
\[
A^2 = A.
\]
(Určete všechny možnosti a dokažte, že žádné jiné neexistují.)
\end{enumerate}

\vfill
\begin{center}
\textit{Konec zadání}
\end{center}

\end{document}
